% from ega0-1-fr.tex
\oldpage[0\textsubscript{I}]{42}
à celle d'un $\mathcal{B}$-homomorphisme
$\mathcal{C}\otimes_{\mathcal{B}}\mathcal{M}\to\mathcal{M}$ vérifiant la
condition d'associativité~; d'où par transport une structure de
$\Psi^*(\mathcal{C})$-Module sur $\Psi^*(\mathcal{M})$.
\end{env}

\begin{env}[4.3.5]
\label{0.4.3.5-fr}
Soit $\mathcal{J}$ un faisceau d'idéaux de $\mathcal{B}$~; comme le foncteur
$\psi^*$ est exact, le $\psi^*(\mathcal{B})$-Module $\psi^*(\mathcal{J})$
s'identifie canoniquement à un faisceau d'idéaux de $\psi^*(\mathcal{B})$~;
l'injection canonique $\psi^*(\mathcal{J})\to\psi^*(\mathcal{B})$ donne alors
un homomorphisme de $\mathcal{A}$-Modules
\[
  \Psi^*(\mathcal{J})
  =\psi^*(\mathcal{J})\otimes_{\psi^*(\mathcal{B})}\mathcal{A}_{[\theta]}
  \longrightarrow\mathcal{A}~;
\]
nous noterons $\Psi^*(\mathcal{J})\mathcal{A}$, ou
$\mathcal{J}\mathcal{A}$ si aucune confusion n'est à craindre, l'image de
$\Psi^*(\mathcal{J})$ par cet homomorphisme. On a donc par définition
\[
  \mathcal{J}\mathcal{A}
  =\theta^\sharp\bigl(\psi^*(\mathcal{J})\bigr)\mathcal{A}
\]
et en particulier, pour tout $x\in X$,
\[
  (\mathcal{J}\mathcal{A})_x
  =\theta_x^\sharp(\mathcal{J}_{\psi(x)})\mathcal{A}_x,
\]
en tenant compte de l'identification canonique des fibres de
$\psi^*(\mathcal{J})$ et de celles de $\mathcal{J}$ (3.7.2). Si
$\mathcal{J}_1$, $\mathcal{J}_2$ sont deux idéaux de $\mathcal{B}$, on a
\[
  (\mathcal{J}_1\mathcal{J}_2)\mathcal{A}
  =\mathcal{J}_1(\mathcal{J}_2\mathcal{A})
  =(\mathcal{J}_1\mathcal{A})(\mathcal{J}_2\mathcal{A}).
\]
Si $\mathcal{F}$ est un $\mathcal{A}$-Module, on posera
$\mathcal{J}\mathcal{F}=(\mathcal{J}\mathcal{A})\mathcal{F}$.
\end{env}

\begin{env}[4.3.6]
\label{0.4.3.6-fr}
Soient $(Z,\mathcal{C})$ un troisième espace annelé,
$\Psi'=(\psi',\theta')$ un morphisme $(Y,\mathcal{B})\to(Z,\mathcal{C})$~;
si $\Psi''$ est le morphisme composé $\Psi'\circ\Psi$, il résulte de la
définition (4.3.1) et de (4.3.3.1) que l'on a
${\Psi''}^*=\Psi^*\circ{\Psi'}^*$.
\end{env}

\subsection{Relations entre images directes et images réciproques}
\label{subsection:0.4.4-fr}

\begin{env}[4.4.1]
\label{0.4.4.1-fr}
Les hypothèses et notations étant celles de (4.2.1), soit $\mathcal{G}$ un
$\mathcal{B}$-Module. Par définition, un homomorphisme
$u:\mathcal{G}\to\Psi_*(\mathcal{F})$ de $\mathcal{B}$-Modules s'appelle
encore un \emph{$\Psi$-morphisme de $\mathcal{G}$ dans $\mathcal{F}$}, ou
simplement un \emph{homomorphisme de $\mathcal{G}$ dans $\mathcal{F}$} et on
l'écrit $u:\mathcal{G}\to\mathcal{F}$ quand aucune confusion n'en résulte.
Se donner un tel homomorphisme revient à se donner, pour tout couple $(U,V)$
où $U$ est un ouvert de $X$, $V$ un ouvert de $Y$ tels que
$\psi(U)\subset V$, un \emph{homomorphisme}
\[
  u_{U,V}:\Gamma(V,\mathcal{G})\longrightarrow\Gamma(U,\mathcal{F})
\]
\emph{de $\Gamma(V,\mathcal{B})$-modules}, $\Gamma(U,\mathcal{F})$ étant
considéré comme $\Gamma(V,\mathcal{B})$-module au moyen de l'homomorphisme
d'anneaux
$\theta_{U,V}:\Gamma(V,\mathcal{B})\to\Gamma(U,\mathcal{A})$~; les $u_{U,V}$
doivent en outre rendre commutatifs les diagrammes (3.5.1.1). Il suffit
d'ailleurs pour définir $u$ de se donner les $u_{U,V}$ lorsque $U$ (resp.
$V$) parcourt une base $\mathfrak{B}$ (resp. $\mathfrak{B}'$) de la topologie
de $X$ (resp. $Y$) et de vérifier la commutativité de (3.5.1.1) pour ces
restrictions.
\end{env}

\begin{env}[4.4.2]
\label{0.4.4.2-fr}
Sous les hypothèses de (4.2.1) et (4.2.6), soient $\mathcal{H}$ un
$\mathcal{C}$-Module, $v:\mathcal{H}\to\Psi'_*(\mathcal{G})$ un
$\Psi'$-morphisme~; alors
\[
  w:\mathcal{H}\xrightarrow{v}\Psi'_*(\mathcal{G})
  \xrightarrow{\Psi'_*(u)}\Psi'_*(\Psi_*(\mathcal{F}))
\]
est un $\Psi''$-morphisme que l'on appelle le \emph{composé de $u$ et $v$}.
\end{env}

\begin{env}[4.4.3]
\label{0.4.4.3-fr}
Nous allons maintenant voir que l'on peut définir un \emph{isomorphisme}
canonique de \emph{bifoncteurs} en $\mathcal{F}$ et $\mathcal{G}$
\[
  \operatorname{Hom}_{\mathcal{A}}(\Psi^*(\mathcal{G}),\mathcal{F})
  \xrightarrow{\sim}
  \operatorname{Hom}_{\mathcal{B}}(\mathcal{G},\Psi_*(\mathcal{F}))
  \tag{4.4.3.1}\label{0.4.4.3.1-fr}
\]
que nous désignerons par $v\mapsto v_\theta^\flat$ (ou simplement
$v\mapsto v^\flat$ si aucune confusion n'est possible)~; nous noterons
$u\mapsto u_\theta^\sharp$, ou $u\mapsto u^\sharp$ l'isomorphisme réciproque.
Cette définition est la suivante~: en composant
$v:\Psi^*(\mathcal{G})\to\mathcal{F}$ avec l'application canonique
$\psi^*(\mathcal{G})\to\Psi^*(\mathcal{G})$, on obtient un homomorphisme de
faisceaux de groupes $v':\psi^*(\mathcal{G})\to\mathcal{F}$, qui est aussi un
homomorphisme de $\psi^*(\mathcal{B})$-modules. On en déduit (3.7.1) un
homomorphisme
${v'}^\flat:\mathcal{G}\to\psi_*(\mathcal{F})=\Psi_*(\mathcal{F})$, qui est
aussi un homomorphisme de $\mathcal{B}$-Modules comme on
